% !TEX root =./ECE444_Practice_main.tex
%
% L2 practice set (Basic Properties and Terminology), ported from the Jupyter
% Book practice.md / practice-solutions.md into the exam class for Gradescope.
% Blank copy: \begin{solution}[h] reserves h of work space. SOLUTIONS copy
% (\iskey=1) prints the red worked solutions.

\fullwidth{\textbf{Learning Objective 1.2: } I can define and calculate
fundamental antenna properties --- gain, directivity, effective aperture,
beamwidth, and sidelobe level --- apply the reciprocity principle that links an
antenna's transmit and receive behavior, and use the Friis transmission
equation to predict received power in a link.
\\[4pt]\normalfont\small Show your work. Write a \textbf{Documentation} line at the
end (\textbf{None} if you did not collaborate).}

\begin{questions}

%%%%%%%%%%%%%%%%%%%%%%%%%%%%%%%%%%%%%%%%%
\question \textbf{Wave equation sanity check.} The 1-D wave equation from a
lossless transmission line is
\eq{ \frac{\partial^{2} v}{\partial z^{2}} = LC \frac{\partial^{2} v}{\partial t^{2}}. }
\begin{parts}

	\part Show by direct substitution that $v(z,t) = f(z - ut)$ is a solution for
	any twice-differentiable $f$, and find $u$ in terms of $L$ and $C$.
		\begin{solution}[1.5in]
		Let $\xi = z - ut$. By the chain rule $\dfrac{\partial^{2} v}{\partial z^{2}} = f^{\prime\prime}(\xi)$
		and $\dfrac{\partial^{2} v}{\partial t^{2}} = u^{2} f^{\prime\prime}(\xi)$. Substituting,
		\eq{ f^{\prime\prime}(\xi) = LCu^{2} f^{\prime\prime}(\xi) \Longrightarrow u = \dfrac{1}{\sqrt{LC}}. }
		\end{solution}

	\begin{minipage}[t]{0.58\linewidth}
		\part For a coaxial cable with $L = 250\ \text{nH/m}$ and $C = 100\ \text{pF/m}$,
		compute the propagation speed $u$ and the velocity factor $u/c$.
		\begin{solution}[1.3in]
		\eq{ LC &= (2.5\times10^{-7})(1.0\times10^{-10}) = 2.5\times10^{-17}\ \text{s}^{2}/\text{m}^{2} \\
			u &= \dfrac{1}{\sqrt{LC}} = 2.0\times10^{8}\ \text{m/s}, \quad \dfrac{u}{c} = \dfrac{2.0}{3.0} \approx 0.67 }
		\end{solution}
	\end{minipage}\hfill%
	\begin{minipage}[t]{0.37\linewidth}
		\raggedleft \vspace{12pt}
		$u$ = \ansbox[1.3in]{$2.0\times10^{8}\ \text{m/s}$}

		\vspace{4pt}
		$u/c$ = \ansbox[1.3in]{$\approx 0.67$}
	\end{minipage}

	\part In one sentence, explain what a ``velocity factor of 0.66'' (typical of
	RG-58 coax) means physically.
		\begin{solution}[0.9in]
		Signals travel along the line at about 66\% of the speed of light in vacuum:
		$u = 1/\sqrt{LC} \approx 0.66c$, set by the line's own $L$ and $C$. Since
		$u = f\lambda$ still holds on the line, a wave at a fixed frequency has a
		wavelength on the line that is $0.66\times$ its free-space value.
		\end{solution}

\end{parts}

\newpage
%%%%%%%%%%%%%%%%%%%%%%%%%%%%%%%%%%%%%%%%%
\question \textbf{Time and space frequencies.} A plane wave in free space is
$\vec{E}(z,t) = \xhat E_{0}\cosp{\omega t - kz}$.
\begin{parts}

	\begin{minipage}[t]{0.58\linewidth}
		\part For a signal at $f = 2.40\ghz$ (Wi-Fi channel 6), compute $\omega$,
		$\lambda$, and $k$.
		\begin{solution}[1.5in]
		\eq{ \omega &= 2\pi f \approx 1.51\times10^{10}\ \text{rad/s} \\
			\lambda &= \dfrac{c}{f} = \dfrac{3.00\times10^{8}}{2.40\times10^{9}} = 0.125\m \\
			k &= \dfrac{2\pi}{\lambda} \approx 50.3\ \text{rad/m} }
		\end{solution}
	\end{minipage}\hfill%
	\begin{minipage}[t]{0.37\linewidth}
		\raggedleft \vspace{12pt}
		$\omega$ = \ansbox[1.5in]{$1.51\times10^{10}\ \text{rad/s}$}

		\vspace{4pt}
		$\lambda$ = \ansbox[1.3in]{$0.125\m$}

		\vspace{4pt}
		$k$ = \ansbox[1.3in]{$\approx 50.3$ rad/m}
	\end{minipage}

	\begin{minipage}[t]{0.58\linewidth}
		\part How far does the wave crest travel in \textbf{one nanosecond}? Express
		your answer in wavelengths.
		\begin{solution}[1in]
		\eq{ c\Delta t = (3.00\times10^{8})(10^{-9}) = 0.30\m = \dfrac{0.30}{0.125} = 2.4\lambda }
		\end{solution}
	\end{minipage}\hfill%
	\begin{minipage}[t]{0.37\linewidth}
		\raggedleft \vspace{12pt}
		\ansbox[1.3in]{$2.4\lambda$}
	\end{minipage}

	\begin{minipage}[t]{0.58\linewidth}
		\part At a fixed point in space, how many full oscillations does the field
		complete in \textbf{one nanosecond}?
		\begin{solution}[0.9in]
		\eq{ f\Delta t = (2.40\times10^{9})(10^{-9}) = 2.4\ \text{cycles} }
		\end{solution}
	\end{minipage}\hfill%
	\begin{minipage}[t]{0.37\linewidth}
		\raggedleft \vspace{12pt}
		\ansbox[1.3in]{$2.4$ cycles}
	\end{minipage}

	\part Sketch (by hand) $\vec{E}$ vs.\ $z$ at $t=0$ over a range of two
	wavelengths, then $\vec{E}$ vs.\ $t$ at $z=0$ over a range of two periods.
	Label the axes.
		\begin{solution}[2in]
		Two identical cosines: one vs.\ $z$ with period $\lambda = 12.5$ cm, one vs.\
		$t$ with period $T = 1/f \approx 417$ ps. The answers to (b) and (c) are the
		\emph{same number} --- that is $\omega/k = c$ in disguise: cycles per meter of
		travel equals cycles per second of elapsed time, related by the speed of light.
		\end{solution}

\end{parts}

\newpage
%%%%%%%%%%%%%%%%%%%%%%%%%%%%%%%%%%%%%%%%%
\question \textbf{Directivity from beamwidths.} A ground-based tracking antenna
has half-power beamwidths $\theta_{1} = 2.5\mydeg$ in azimuth and
$\theta_{2} = 1.8\mydeg$ in elevation.
\begin{parts}

	\begin{minipage}[t]{0.58\linewidth}
		\part Estimate the directivity in dBi using the pencil-beam approximation
		$D \approx 41{,}253/(\theta_{1}^{\circ}\theta_{2}^{\circ})$.
		\begin{solution}[1.2in]
		\eq{ D &\approx \dfrac{41{,}253}{(2.5)(1.8)} = \dfrac{41{,}253}{4.5} \approx 9170 \\
			D_\text{dBi} &= 10\mylog{9170} \approx 39.6\dbi }
		\end{solution}
	\end{minipage}\hfill%
	\begin{minipage}[t]{0.37\linewidth}
		\raggedleft \vspace{12pt}
		$D$ = \ansbox[1.3in]{$\approx 39.6\dbi$}
	\end{minipage}

	\begin{minipage}[t]{0.58\linewidth}
		\part You measure the antenna's radiation efficiency at $88\%$ and its impedance
		mismatch $\left|\Gamma\right| = 0.25$. What is the realized gain at boresight, in dBi?
		\begin{solution}[1.6in]
		\eq{ G &= \eta_\text{rad}D = 0.88(9170) \approx 8070 \to 39.1\dbi \\
			1 - \left|\Gamma\right|^{2} &= 1 - 0.0625 = 0.9375 \\
			G_\text{re} &= 0.9375(8070) \approx 7560 \to 38.8\dbi }
		\end{solution}
	\end{minipage}\hfill%
	\begin{minipage}[t]{0.37\linewidth}
		\raggedleft \vspace{12pt}
		$G_\text{re}$ = \ansbox[1.3in]{$\approx 38.8\dbi$}
	\end{minipage}

	\part In one sentence, explain the difference between $D$, $G$, and
	$G_\text{re}$ for this antenna.
		\begin{solution}[1.1in]
		$D$ is purely geometric (pattern shape only); $G$ folds in the ohmic and
		dielectric losses inside the antenna; $G_\text{re}$ additionally subtracts the
		impedance-mismatch loss at the terminals --- it is what you would actually measure
		with a source connected through a transmission line.
		\end{solution}

\end{parts}

\newpage
%%%%%%%%%%%%%%%%%%%%%%%%%%%%%%%%%%%%%%%%%
\question \textbf{Effective area.}
\begin{parts}

	\begin{minipage}[t]{0.58\linewidth}
		\part A parabolic dish for a satellite ground terminal has $G = 47\dbi$ at
		$f = 12\ghz$. Compute the effective area $A_{e} = G\lambda^{2}/(4\pi)$.
		\begin{solution}[1.5in]
		\eq{ \lambda &= \dfrac{c}{f} = 0.0250\m, \quad G = 10^{4.7} \approx 50{,}120 \\
			A_{e} &= \dfrac{G\lambda^{2}}{4\pi} = \dfrac{50{,}120(6.25\times10^{-4})}{4\pi} \approx 2.49\ \text{m}^{2} }
		\end{solution}
	\end{minipage}\hfill%
	\begin{minipage}[t]{0.37\linewidth}
		\raggedleft \vspace{12pt}
		$A_{e}$ = \ansbox[1.3in]{$\approx 2.49\ \text{m}^{2}$}
	\end{minipage}

	\begin{minipage}[t]{0.58\linewidth}
		\part The dish is physically $2.4\m$ in diameter. What is its aperture
		efficiency $\eta_\text{ap} = A_{e}/A_\text{phys}$? Is that reasonable?
		\begin{solution}[1.4in]
		\eq{ A_\text{phys} &= \pi(1.2)^{2} \approx 4.52\ \text{m}^{2} \\
			\eta_\text{ap} &= \dfrac{2.49}{4.52} \approx 0.55 }
		Realistic --- real dishes fall in $\eta_\text{ap} \approx 0.5$--$0.7$ once
		illumination taper, spillover, feed/strut blockage, and surface roughness are counted.
		\end{solution}
	\end{minipage}\hfill%
	\begin{minipage}[t]{0.37\linewidth}
		\raggedleft \vspace{12pt}
		$\eta_\text{ap}$ = \ansbox[1.3in]{$\approx 0.55$}
	\end{minipage}

	\begin{minipage}[t]{0.58\linewidth}
		\part Keeping the same physical dish but re-illuminating it efficiently at
		$f = 30\ghz$ (Ka-band), what gain (dBi) would you expect if $\eta_\text{ap}$
		is unchanged?
		\begin{solution}[1.4in]
		$A_{e}$ fixed $\Rightarrow G = 4\pi A_{e}/\lambda^{2} \propto f^{2}$. From 12 to
		30 GHz is $2.5\times$ in $f$, i.e.\ $6.25\times$ in gain $= 10\mylog{6.25} \approx 8.0\db$.
		\eq{ G_{30} \approx 47 + 8 = 55\dbi }
		(Optimistic: higher $f$ exposes surface roughness that erodes $\eta_\text{ap}$.)
		\end{solution}
	\end{minipage}\hfill%
	\begin{minipage}[t]{0.37\linewidth}
		\raggedleft \vspace{12pt}
		$G_{30}$ = \ansbox[1.3in]{$\approx 55\dbi$}
	\end{minipage}

\end{parts}

\newpage
%%%%%%%%%%%%%%%%%%%%%%%%%%%%%%%%%%%%%%%%%
\question \textbf{Reading a real pattern.} Find a published radiation-pattern
plot for a real antenna --- a manufacturer datasheet, application note, or
journal figure with a dB-labeled axis --- and hand in the annotated plot with
your readings. \emph{Cite the source; answers are antenna-dependent, so the
solutions grade method, not a specific number.}
\begin{parts}
	\part Boresight direction (degrees off nose).
		\begin{solution}[0.5in] Marked at the true peak of the main lobe. \end{solution}
	\part Half-power beamwidth (HPBW).
		\begin{solution}[0.5in] Read at the $-3\db$ points, measured consistently on both sides. \end{solution}
	\part First-null beamwidth (FNBW), or state if no null is visible in range.
		\begin{solution}[0.5in] Read at the first nulls; ``no null visible'' is legitimate if the cut is monotonic to the plot edge. \end{solution}
	\part Peak sidelobe level (SLL), in dB below the main lobe.
		\begin{solution}[0.5in] The highest sidelobe peak, in dB below the main lobe. \end{solution}
	\part Front-to-back ratio (F/B), in dB.
		\begin{solution}[0.6in] Main lobe minus the level at $180\mydeg$; ``no distinct back lobe'' is acceptable with justification. The habit here is reading pattern plots \emph{consistently} --- every later lesson reuses these same measurements. \end{solution}
\end{parts}

\newpage
%%%%%%%%%%%%%%%%%%%%%%%%%%%%%%%%%%%%%%%%%
\question \textbf{Friis link with gain and effective area.} A UHF ground station
at $f = 400\mhz$ transmits $P_{t} = 20\watts$ into an antenna with
$G_{t} = 8\dbi$. A cubesat at $r = 800\text{ km}$ carries a receive antenna with
$G_{r} = 3\dbi$ and is on boresight for both antennas. Take $\lambda = c/f = 0.75\m$.
\begin{parts}

	\begin{minipage}[t]{0.58\linewidth}
		\part Compute the on-axis power density $S_\text{inc}$ at the cubesat.
		\begin{solution}[1.5in]
		$G_{t} = 8\dbi \to 6.31$.
		\eq{ S_\text{inc} &= \dfrac{P_{t} G_{t}}{4\pi r^{2}} = \dfrac{20(6.31)}{4\pi(8.0\times10^{5})^{2}} \\
			&\approx 1.57\times10^{-11}\ \text{W/m}^{2} }
		\end{solution}
	\end{minipage}\hfill%
	\begin{minipage}[t]{0.37\linewidth}
		\raggedleft \vspace{12pt}
		$S_\text{inc}$ = \ansbox[1.6in]{$1.57\times10^{-11}\ \text{W/m}^{2}$}
	\end{minipage}

	\begin{minipage}[t]{0.58\linewidth}
		\part Compute the cubesat receive antenna's effective area $A_{e,r}$.
		\begin{solution}[1.2in]
		$G_{r} = 3\dbi \to 2.00$.
		\eq{ A_{e,r} = \dfrac{G_{r}\lambda^{2}}{4\pi} = \dfrac{2.00(0.5625)}{4\pi} \approx 0.0895\ \text{m}^{2} }
		\end{solution}
	\end{minipage}\hfill%
	\begin{minipage}[t]{0.37\linewidth}
		\raggedleft \vspace{12pt}
		$A_{e,r}$ = \ansbox[1.3in]{$\approx 0.0895\ \text{m}^{2}$}
	\end{minipage}

	\begin{minipage}[t]{0.58\linewidth}
		\part Compute the received power $P_{r} = S_\text{inc}A_{e,r}$ and express
		the result in dBm.
		\begin{solution}[1.3in]
		\eq{ P_{r} &= (1.57\times10^{-11})(0.0895) \approx 1.40\times10^{-12}\watts \\
			P_{r,\text{dBm}} &= 10\mylog{\dfrac{1.40\times10^{-12}}{10^{-3}}} \approx -88.5\dbm }
		\end{solution}
	\end{minipage}\hfill%
	\begin{minipage}[t]{0.37\linewidth}
		\raggedleft \vspace{12pt}
		$P_{r}$ = \ansbox[1.3in]{$\approx -88.5\dbm$}
	\end{minipage}

	\part The cubesat receiver needs $-82.5\dbm$ at its terminals to close the link
	with the margin the program office demands, so this design is $6\db$ short. Friis,
	$P_{r} = P_{t}G_{t}G_{r}\lp \dfrac{\lambda}{4\pi r} \rp^{2}$, says you can buy those
	$6\db$ from $P_{t}$, from $G_{t}$, or from $G_{r}$. Pick one and defend it, saying
	what each of the three costs.
		\begin{solution}[2.0in]
		All three enter Friis as a plain multiplication, so a decibel is a decibel no
		matter which factor supplies it --- the choice is an engineering one, not a
		physics one.
		\textbf{$P_{t}$:} $6\db$ means going from $20\watts$ to $80\watts$. That is the
		easiest change to make on the ground, but it costs prime power, cooling, and
		amplifier linearity, and on a licensed band it may run straight into an EIRP
		limit that more power cannot buy its way past.
		\textbf{$G_{t}$:} $8\dbi \rightarrow 14\dbi$ roughly halves the ground antenna's
		beamwidth. Ground real estate and structure are cheap, so this is usually the
		right answer --- the price is that the station must now track the pass accurately,
		and EIRP rises exactly as it would with more power.
		\textbf{$G_{r}$:} $6\db$ more on the cubesat means a larger deployable antenna
		with a narrower beam, on a spacecraft whose attitude control may not be able to
		point it. Mass and pointing are the scarcest resources in the whole system.
		\textbf{Recommendation: buy it at $G_{t}$.} It is the cheapest decibel in the
		link and it sits on the end where size, power, and pointing hardware are all
		affordable --- provided the EIRP limit leaves room.
		\end{solution}

\end{parts}

\vspace{1.5em}
\noindent\textbf{Documentation:} \rule[-2pt]{0.6\linewidth}{0.4pt}

\end{questions}
